\subsection{Анализ существующих решений для построения графов знаний}

В настоящее время на рынке представлено множество решений
для управления знаниями, которые автоматизированное строят ГЗ.
Они варьируются от полноценных экосистем, предназначенных для корпоративного использования
до узконаправленных инструментов, решающих конкретные задачи.
Проведенный анализ позволил разделить решения на три группы, выявить ограничения и хорошие практики
существующих решений, а также сформировать базовый набор технологий,
применяемый при построении ГЗ в контексте систем управления знаниями.

\textbf{Готовые решения для корпоративного использования.}
В данную категорию включены полноценные экосистемы,
позволяющие автоматически трансформировать пользовательские текстовые данные
в графы знаний при помощью LLM.
Представителями этой категории являются такие решения как
Neo4j Graph Builder~\cite{neo4j_graph_builder} и Memgraph Unstructured2Graph~\cite{memgraph_u2g}.
Сами Neo4j и Memgraph являются графовыми базами данных,
предоставляющими широкий функционал «из коробки»,
в том числе возможность извлечения сущностей и связей из
пользовательских текстовых данных различных популярных форматов
для построения ГЗ на их основе.

Поскольку данная работа фокусируется на построении ГЗ,
в этих решениях будет изучен исключительно ответственный за это модуль.
Оба решения имеют открытый исходный код этих модулей.
За основу взят пайплайн Neo4j, однако логика Memgraph крайне схожа.
Изученный конвейер состоит из нескольких этапов:

\begin{enumerate}
      \item Парсинг и чанкинг. Исходные данные разбиваются на фрагменты размером
            4000 символов и перекрытием 200 символов [https://github.com/neo4j-labs/llm-graph-builder].
      \item Извлечение сущностей. Каждый чанк передается в LLM с инструкциями для извлечения
            троек \texttt{(Entity A — [relation type] — Entity B)}.
      \item Разрешение кореференций. Извлеченные имена сущностей
            дедуплицируются с помощью большой языковой модели
            для избежания конфликтов
            (например, сущности «ИИ» и «Искусственный интеллект» объединяются в одну с именем
            «Искусственный интеллект»).
      \item Запись данных. Обработанные данные записываются в базу данных.
            Они формируют два графа:
            лексический, сохраняющий последовательность фрагментов в файле, и семантический,
            содержащий связи между сущностями.
      \item Создание сообществ. Используются
            графовые алгоритмы (Louvain, PageRank) для поиска сообществ,
            что позволяет создавать кластерные узлы (хабы),
            увеличивающие эффективность навигации RAG.
\end{enumerate}

Neo4j и Memgraph используют LangChain~\cite{Chase_langchain} для работы с LLM,
что позволяет использовать множество провайдеров,
а также запущенные локально модели.

Данные системы обладают рядом преимуществ:
\begin{enumerate}
      \item Поддержка множества форматов текстовых данных,
            таких как PDF, HTML и DOCX.
      \item Построение двухуровневых графов: лексический,
            отражающий последовательность фрагментов текстов,
            и граф сущностей, соединенных типизированными связями.
            Это упрощает сбор всего содержимого одной заметки, например,
            в случае необходимости пересказа ее содержимого
      \item Формирование сообществ с описаниями,
            что экономит контекст при ответе на общие вопросы,
            затрагивающие множество заметок на общую тему.
\end{enumerate}

Однако они также имеют ряд критических недостатков, ставящих
под сомнение использование подобных решений для персональных баз данных.
\begin{enumerate}
      \item Отсутствие поддержки частичного обновления. Даже при минимальных модификациях
            файлов необходимо полностью перезапускать дорогостоящий пайплайн для построения ГЗ.
      \item Разрыв между данными и графом. Поскольку в этих сервисах отсутствуют
            механизмы синхронизации содержимого заметок, пользователю необходимо
            вручную загружать обновленные файлы на сервер базы данных,
            что увеличивает вероятность ошибок
            и делает процесс более трудоемким.
      \item Инфраструктурные издержки. Для обработки данных и взаимодействия с ГЗ необходима
            работающая база данных, что ведет к потреблению дополнительных вычислительных ресурсов,
            и, как следствие, дополнительным затратам, если сервис работает на выделенном сервере.
      \item Чрезмерное доверие к LLM. В пайплайнах этих решений отсутствует фильтрация кандидатов.
            В них вся работа по извлечению сущностей и анализу связей
            полностью лежит на больших языковых моделях, которым свойственны «галлюцинации».
            Это влечет создание бессмысленных и выдуманных связей,
            которые делают извлечение информации неэффективным.
\end{enumerate}

\textbf{Академические автоматизированные пайплайны.}
В качестве альтернативы решениям, опирающимся на
использование тяжеловесных LLM было найдено решение~\cite{Ain_2025},
эффективно использующее сравнительно небольшие модели и другие
методики для построения ГЗ.

Его пайплайн состоит их следующих этапов:
\begin{enumerate}
      \item Чанкинг. Файлы разбиваются на фрагменты, каждый из которых будет
            обрабатываться независимо, то есть для каждого чанка будет создан независимый ГЗ.
      \item Извлечение сущностей. Вместо больших языковый моделей данное решение использует
            компактную модель SqueezeBERT~\cite{Iandola_2020}. Процесс извлечения аналогичный:
            в модель передается фрагмент текста и инструкции.
      \item Привязка к онтологиям. Для стандартизации имен сущностей и снижения галлюцинаций
            все извлеченные сущности сверяются с базами данных, такими как DBpedia или Wikidata.
      \item Вычисление семантического сходства. Для создания связей между сущностями
            используется сравнение косинусного сходства их плотных эмбеддингов,
            созданных моделью Sentence-BERT~\cite{Reimers_2019}, с пороговым значением.
      \item Слияние графов. Все ГЗ, построенные для фрагментов
            текстов объединяются в единую сеть на основе общих узлов и связей.
\end{enumerate}

Ключевые достоинства этого подхода включают:
\begin{enumerate}
      \item Вычислительную эффективность.
            Использование небольших моделей машинного обучения сильно экономит
            вычислительные ресурсы и позволяет использовать решение полностью локально.
      \item Использование онтологий. Это решение упрощает этап разрешения кореференций,
            уменьшая вероятность галлюцинаций
            и убирая необходимость использования LLM для дедупликации.
\end{enumerate}

У данного решения есть несколько недостатков,
делающими этот подход неподходящим для решения главной проблемы этой работы.

\begin{enumerate}
      \item Обязательное участие человека. Поскольку авторам важно качество полученного ГЗ,
            сгенерированный граф тщательно корректируется человеком, что делает процесс
            лишь частично автоматизированным.
      \item Малое разнообразие сущностей. Использование жестких списков допустимых сущностей
            может повлечь ошибочную фильтрацию непопулярных терминов или аббревиатур.
\end{enumerate}

\textbf{Облачные приложения для заметок.}
Индустрия также имеет примеры приложений для ведения заметок,
которым удалось интегрировать ГЗ для увеличения эффективности их использования.

Облачные сервисы, такие как Mem.ai~\cite{Mem_ai}, являются
программным обеспечением как услугой (SaaS), полностью разгружая пользователя
от организации и связывания заметок.
Mem.ai, по аналогии с MemGraph и Neo4j, строит двухуровневый граф
(лексический слой,
сохраняющий последовательность фрагментов в файле, и семантический,
содержащий связи между сущностями).
Это обеспечивает эффективную навигацию алгоритмам RAG.
В данном сервисе ГЗ используются как для подсказок во время написания заметок,
так и для извлечения и синтеза информации посредством интегрированного чат-бота.

Несмотря на все удобства, у Mem.ai существуют недочеты,
касающиеся приватности и обработки данных:
\begin{enumerate}
      \item Закрытый исходный код. Пользователи не имеют доступа к логике построения ГЗ
            и обработки их заметок.
      \item Привязка к поставщику. Данные хранятся на серверах компании,
            что ставит под сомнение их безопасность,
            и делает их экспорт в случае закрытия сервиса затруднительным:
            сервис предоставляет возможность экспорта сырого текста заметок,
            но не ГЗ.
\end{enumerate}

\textbf{Плагины Obsidian.}
Приложение имеет поддержку пользовательских расширений,
которые добавляют новый функционал в программу.
По результатам исследования выяснилось,
что расширений, позволяющих автоматизировать процесс построения ГЗ не существует.
Несмотря на это, есть плагины-ассистенты, упрощающие некоторые его этапы.

\begin{enumerate}
      \item \textbf{Smart Connections}~\cite{smart_connections}. Плагин, использующих векторный поиск для
            формирования списка схожих заметок.
            Инструмент способен находить только семантически близкие заметки.
            Кроме того, плагин не проводит никакой дополнительной фильтрации результатов,
            а также не типизирует результаты, чем перекладывает
            большую часть работу на пользователя.
      \item \textbf{Graph Analysis}~\cite{graph_analysis}. Инструмент, подсказывающий связи на основе существующих.
            Данное решение опирается только на уже построенный граф для вычисления статистик,
            на основе которых предлагает пользователю новые связи или хабы.
            Данный плагин может хорошо работать
            как дополнение к автоматизированному пайплайну построения ГЗ,
            но не как его замена.
\end{enumerate}

\textbf{Выводы.}
Решения, представленные на рынке в данный момент, заставляют пользователей идти на компромиссы.
С одной стороны есть либо полноценные дорогостоящие и требовательные к ресурсам базы данных,
либо закрытые облачные решения для ведения заметок,
которые все же решают задачу автоматизированного построения ГЗ.
С другой стороны существуют открытые легковесные инструменты,
позволяющие упростить некоторые этапы построения,
но требующие вмешательства человека.

Решение с открытым исходным кодом, способное вычислительно эффективно находить,
фильтровать, типизировать и сохранять связи, а также поддерживающее работу полностью локально,
имело бы огромное преимущество на фоне изученных решений.
Плюсом будет возможность модификации построенного графа пользователем,
как и поддержка инкрементальных обновлений.
